%! Equivalent Representations of Polynomial Expressions — Unit 2, Lesson 2.6 — Student Packet Cover Sheet
\documentclass[10pt]{article}
\usepackage{precalculus-article}
\usepackage{precalculus-boxes}
\usepackage{ltablex}
\keepXColumns

\begin{document}

% Full-bleed plum banner
\begin{tikzpicture}[remember picture, overlay]
  \fill[plum] (current page.north west)
    rectangle ([yshift=-0.9in]current page.north east);
\end{tikzpicture}
\vspace{-0.8in}

\begin{center}
  {\color{white}\textbf{\LARGE Precalculus}}\\[4pt]
  {\color{lilacmid}\large\textbf{Unit 2: Polynomial Functions}}\\[6pt]
  {\color{white}\textbf{\large Lesson 2.6 \quad Equivalent Representations of Polynomial Expressions}}
\end{center}

\vspace{0.22in}
\namedateperiod
\vspace{0.18in}

\begin{learningtargetbox}
\small
\begin{enumerate}[leftmargin=1.5em, itemsep=5pt]
  \item I can decide whether the factored form or the standard form of a
        polynomial is the useful one for a given question, and I can explain
        what each form hands me --- real zeros and $x$-intercepts from one,
        leading term, end behavior, and $y$-intercept from the other.
  \item I can rewrite a polynomial in factored form as a polynomial in
        standard form by multiplying the factors out.
  \item I can use one row of Pascal's Triangle and the binomial theorem to
        expand $(x+c)^n$ without multiplying $n$ copies of the binomial.
  \item I can divide one polynomial by another using long division and write
        the result as $f(x) = g(x)q(x) + r(x)$, where the remainder has lower
        degree than the divisor.
  \item I can use a remainder of $0$ to confirm that a divisor is a factor,
        and I can use both forms of the same polynomial to answer questions
        about a situation.
\end{enumerate}
\end{learningtargetbox}

\vspace{0.14in}

\begin{tocbox}
\small
\renewcommand{\arraystretch}{1.5}
\begin{tabularx}{\linewidth}{@{} c l X r @{}}
\rowcolor{lilac}
\textbf{\#} & \textbf{Component} & \textbf{Description} & \textbf{Score} \\
\midrule
1 & Warm-Up        & Spiral review of prerequisite skills                     & \blank{1.2cm} \\
2 & Guided Notes   & Two forms, the binomial theorem, and long division       & \blank{1.2cm} \\
3 & Group Activity & Same Polynomial, Two Jobs: choosing and converting forms & \blank{1.2cm} \\
4 & Exit Ticket    & Independent check on forms, expanding, and dividing      & \blank{1.2cm} \\
5 & Homework       & Practice and extension                                   & \blank{1.2cm} \\
\midrule
  & \multicolumn{2}{r}{\textbf{Total}} & \blank{1.2cm} \\
\end{tabularx}
\end{tocbox}

\vspace{0.14in}

\begin{remindbox}
\small
Nothing today changes \textit{which} polynomial you are looking at --- only how
it is written. The factored form and the standard form are the same function
in two costumes, and each one answers questions the other cannot. The skill is
knowing which costume the question wants, and being able to change into it.
\end{remindbox}

\end{document}
